%% one_page_summary.tex
%% LaTeX fragment; \input at the start of \section{Introduction}.
%% Requires: booktabs, amsthm, amsmath (all loaded in the master file).

%%---------------------------------------------------------------------------
\paragraph{Framework.}
An \emph{explanation system} $\mathcal{S} = (\Params, \sH, Y, \mathsf{obs},
\mathsf{exp}, \incomp)$ pairs a parameter space $\Params$ with a hypothesis
class $\sH$ equipped with a native explanation map $\mathsf{obs}: \sH \to Y$
and a practitioner explanation map $\mathsf{exp}: \sH \to Y$, where $\incomp$
is an incompatibility relation on $Y$.
A system has the \emph{Rashomon property} if there exist observationally
equivalent hypotheses $h \sim h'$ (same $\theta \in \Params$) such that
$\mathsf{obs}(h) \incomp \mathsf{obs}(h')$---that is, the data do not
determine the explanation even after fixing the task.
Three desiderata for $\mathsf{exp}$ are then natural: \emph{faithfulness}
($\mathsf{exp}$ never contradicts $\mathsf{obs}$), \emph{stability} ($h \sim
h'$ implies $\mathsf{exp}(h) = \mathsf{exp}(h')$), and \emph{decisiveness}
($\mathsf{exp}$ inherits every incompatibility pair from $\mathsf{obs}$).

%%---------------------------------------------------------------------------
\begin{theorem}[Explanation Impossibility]
\label{thm:universal-summary}
No explanation $\mathsf{exp}$ can simultaneously be faithful, stable, and
decisive for any system with the Rashomon property.
\end{theorem}

\noindent\emph{Proof sketch.}
Let $h \sim h'$ with $\mathsf{obs}(h) \incomp \mathsf{obs}(h')$.
Decisiveness forces $\mathsf{exp}(h) \incomp \mathsf{exp}(h')$.
Stability forces $\mathsf{exp}(h) = \mathsf{exp}(h')$.
No explanation can be both incompatible with itself and equal to itself;
contradiction.
(Lean-verified in \texttt{ExplanationSystem.lean} with zero axiom
dependencies.)\hfill$\square$

%%---------------------------------------------------------------------------
\paragraph{Tightness.}
The impossibility is tight: each of the three pairs of properties is
simultaneously achievable, confirming all three conditions are necessary.
Dropping faithfulness yields fixed-rule explanations (stable and decisive
but ignoring model internals).
Dropping stability yields single-model explanations such as \SHAP{} or
GradCAM (faithful and decisive but varying on retrain).
Dropping decisiveness yields $G$-invariant aggregation methods such as
\DASH{}, CPDAGs, and ensemble probes (faithful in expectation and stable but
returning equivalence classes rather than single answers).

%%---------------------------------------------------------------------------
\paragraph{Nine instances.}
Table~\ref{tab:summary-instances} records the nine explanation types covered
by the framework, their Rashomon mechanism, and the empirical instability
magnitude measured across functionally equivalent models.

\begin{table}[h]
\centering
\small
\caption{Nine instances of the Explanation Impossibility with
empirical instability results.}
\label{tab:summary-instances}
\begin{tabular}{llll}
\toprule
Instance & Rashomon mechanism & Instability metric & Result \\
\midrule
Attribution   & Collinearity / Rashomon set & Feature-rank flip rate
              & $>50\%$ (GBDT)\\
Attention     & Weight-space multiplicity   & Argmax flip rate
              & $60.0\%$ \\
Counterfactual & Model multiplicity         & Direction flip rate
              & $23.5\%$ \\
Concept       & Hidden-layer symmetry       & Direction instability
              & $0.90$ (cosine) \\
Causal        & Markov equivalence class    & DAG flip rate
              & $100\%$ (within MEC) \\
Model Sel.    & Rashomon multiplicity       & Best-model flip rate
              & $80.0\%$ \\
GradCAM       & Weight perturbation         & Peak-pixel flip rate
              & $9.6\%$ \\
LLM Expl.     & Attention multiplicity      & Citation flip rate
              & $34.5\%$ \\
Mech.\ Interp. & Circuit multiplicity       & Circuit flip rate
              & (argued) \\
\bottomrule
\end{tabular}
\end{table}

%%---------------------------------------------------------------------------
\paragraph{Resolution and implications.}
Every instance admits a $G$-invariant resolution: aggregating $\mathsf{exp}$
over the orbit of equivalent hypotheses restores faithfulness-in-expectation
and stability at the cost of decisiveness, and \DASH{}, CPDAGs, and ensemble
concept probes are concrete realizations.
The impossibility is ubiquitous---any learning problem whose parameter
dimension exceeds the data dimension admits a Rashomon set of positive
measure---so the trilemma applies to virtually all modern overparameterized
models, and practitioners must explicitly choose which property to sacrifice.
